% =============================================================================
\section{Divisional charts (vargas)}
\label{sec:vargas}
% =============================================================================

\subsection{What PL7 offers}

The manual lists sixteen Parashari vargas (D1, D2, D3, D4, D7, D9, D10, D12, D16, D20,
D24, D27, D30, D40, D45, D60) and four Tajika vargas: Panchamsha D5, Shashtamsha D6,
Ashtamsha D8 and Ekadashamsha D11 \tagDOC\ \parencite[pp.~40, 52--53]{geovision2006pl7}. It
also mentions ``alternate'' vargas marked II or III, without defining them. PL7's own
documentation gives \emph{no} computational rule for any varga; the bundled translation of
BPHS (\texttt{CHAP/CHAP06.TXT}) states the classical rules for D9, D10 and D24 only.

\subsection{What the code actually does}

All vargas go through one dispatcher, \texttt{fcn.0057d9cf}, a 31-way \texttt{switch} on
the varga's position in the list. A byte table at \texttt{0x5f47b8} gives the divisor of
each position: $1, 2, 3, 4, 7, 9, 10, 12, 16, 20, 24, 27, 30, 40, 45, 60$ (positions
1--16), then $5, 6, 8, 11$ (the Tajika four), then $150$ (Nadiamsha), then ten alternate
vargas with divisors $2, 2, 3, 4, 10, 12, 30, 5, 16, 11$ \tagBIN. There are two kinds of
routine:

\begin{itemize}[leftmargin=*]
  \item \textbf{Cyclic.} \texttt{fcn.0057d1e1} returns $n\lambda \bmod 360^\circ$, where
  $\lambda$ is the sidereal longitude and $n$ the divisor. Counting continuously from Aries
  in this way \emph{is} the standard Parashari rule for D7, D9, D16, D20 and D27; for D9 it
  reproduces ``movable from itself, fixed from the 9th, dual from the 5th''.
  \item \textbf{Sign-dependent.} D3, D4, D10, D12, D24, D30, D40, D45, D60, D5, D6, D8,
  D11 and two alternates have their own routines, which take the sign number $s$ (Aries
  $=1$) and the degrees $d$ within the sign.
\end{itemize}

I re-implemented the routines you asked about exactly as decoded
(\texttt{scripts/pl7\_vargas.py}); \autoref{tab:vargas} gives each rule and the starting
sign for every sign.

\begin{table}[htbp]
  \centering
  \small
  \caption[Decoded varga rules]{Varga rules decoded from \texttt{PL7.exe} \tagBIN.
  ``Start'' is the sign of the first amsa for Aries, Taurus, Gemini \ldots\ Pisces, taken
  from the re-implementation. $s$ = sign number (Aries $=1$); $d$ = degrees in the sign.}
  \label{tab:vargas}
  \begin{tabularx}{\textwidth}{>{\raggedright\arraybackslash}p{2.5cm} l >{\raggedright\arraybackslash}X >{\raggedright\arraybackslash}p{3.7cm}}
    \toprule
    Varga & Routine & Rule & Start (Ar \ldots Pi) \\
    \midrule
    D9 Navamsha & \texttt{0057d1e1} & $9\lambda \bmod 360$ (cyclic from Aries) & Ar Cp Li Cn, repeating \\
    D10 Dashamsha & \texttt{0057d55c} & odd $s$: from the sign itself; even $s$: from the 9th & Ar Cp Ge Pi Le Ta Li Cn Sg Vi Aq Sc \\
    D24 Chatur\-vimshamsha & \texttt{0057d675} & odd $s$: from Leo; even $s$: from Cancer & Le Cn, repeating \\
    D6 Shashtamsha & \texttt{0057d4dc} & odd $s$: from Aries; even $s$: from Libra & Ar Li, repeating \\
    D8 Ashtamsha & \texttt{0057d510} & movable: from Aries; fixed: from Sagittarius; dual: from Leo & Ar Sg Le, repeating \\
    D11 Ekadashamsha & \texttt{0057d5a6} & $(11d/30 + 1 - s)\times 30$, identical to $11\lambda \bmod 360$ (cyclic from Aries) & Ar Pi Aq Cp Sg Sc Li Vi Le Cn Ge Ta \\
    D11 alternate (position 31) & \texttt{0057d5d6} & $(1 - 11d/30)\times 30$: every sign starts at Aries and counts \emph{backwards} & Ar for all (2nd amsa Pi) \\
    D16 alternate (position 30) & \texttt{0057d5fb} & odd $s$: forwards from the sign; even $s$: backwards & Ar Le Ge Li Le Sg \ldots \\
    \bottomrule
  \end{tabularx}
\end{table}

\begin{notebox}\ksimple On D11: PL7 has two Ekadashamsha rules and no rule named
Rudramsha or ``Iyer''. Its main D11 is the plain continuous count from Aries; its alternate
D11 counts backwards from Aries for every sign. If the ``Iyer'' D11 you have in mind is a
third rule, PL7 does not implement it, and Kala Chakra has no D11 at all. Given the rule,
it can be added to Kala Chakra as a single function.\end{notebox}

Varga strength (Vimshopaka) is offered on 6-, 7-, 10- and 16-varga schemes with scores of
5--20, and the planetary table's ``vimshopaka bala'' column uses the 10-varga scheme
\tagDOC\ \parencite[p.~145]{geovision2006pl7}. The per-varga weights are not documented and
were not decoded.

% =============================================================================
\section{Bhava chalit (the Bhava chart)}
\label{sec:bhava}
% =============================================================================

PL7 does not use the name ``Bhava Chalit''; the equivalent is the \emph{Bhava chart}
\tagDOC. Three house divisions are offered as special charts
\parencite[p.~56]{geovision2006pl7}:

\begin{enumerate}[leftmargin=*]
  \item \textbf{Bhava (Sripati)}, centred on the ascendant. The tutorial
  (\texttt{Extras1/pib11770}) defines it: the ascendant and MC are the madhyas (centres) of
  houses 1 and 10 and the opposite points those of 7 and 4; ``the spaces in between are
  further divided in 3 equal portions''; and ``the Bhava Sandhis are placed halfway between
  each pair of adjacent houses''. This is Porphyry-style trisection of each quadrant,
  applied to bhava \emph{centres}.
  \item \textbf{Bhava (equal)}: each house spans $15^\circ$ either side of a madhya set by
  the ascendant degree.
  \item \textbf{Krishnamurti chart (Placidus)}, in which the ascendant is the
  \emph{beginning} of the first house; this is the KP cusp chart of \autoref{sec:placidus}.
\end{enumerate}

\begin{keybox}\kidea PL7's ``bhava chalit'' is a Sripati chart: \emph{centres} at Asc and MC
with quadrant trisection. Kala Chakra's houses are KP-style \emph{beginnings}
(Topocentric cusps). The two answer different questions, so a planet can legitimately sit
in different houses in the two programs even with identical longitudes.\end{keybox}

% =============================================================================
\section{Yogi, Avayogi and Duplicate Yogi}
\label{sec:yogidoc}
% =============================================================================

The manual never states the Yogi formula; the binary does (\autoref{sec:yogi}): the Yogi
point is Sun $+$ Moon $+$ $93^\circ 20'$ \tagBIN. The report templates print ``Yogi
Point'', ``Yogi'' (the lord of the Yogi point's nakshatra) and a combined field
``Avayogi/Duplicate Yogi'' (\texttt{Extras1/pib337}) \tagDOC. The dictionary defines
Avayogi only as the point opposite the Yogi point \parencite[p.~208]{geovision2006pl7};
the Duplicate Yogi is not defined at all. Definitions of both vary between authors, and PL7's documentation does not say which it
uses; I could not settle them from the code within the time available, so they are left
open here rather than filled in from general knowledge. Kala Chakra computes
none of these points.

% =============================================================================
\section{Dashas}
\label{sec:dashas}
% =============================================================================

\subsection{Systems and levels}

PL7 offers at least twenty dasha systems (Vimshottari, Ashtottari, Shodashottari,
Dwadashottari, Panchottari, Shatabdika, Chaturashitisama, Dwisaptatisama, Shastihayani,
Shattrimshatsama, Kalachakra, Chara, Sthira, Niryana Shoola, Drig, Yogini, Navamsha,
Tribhagi, Tribhagi~40, Narayana, and more below the visible part of the list), shown to at
most five levels: maha, antar, pratyantar, sookshma and prana \tagDOC\
\parencite[pp.~42--44]{geovision2006pl7}. The reference body can be the Moon (default), the
ascendant or any planet. The balance at birth is computed ``by degree'' (fraction of
$13^\circ 20'$ traversed, the default) or ``by time''. Kala Chakra computes Vimshottari only,
but to \emph{six} levels (adding deha), by degree from the Moon, anchored once at the chart
moment.

\subsection{The dasha year: a hidden discrepancy}

The Dashas tab offers ``360 days or 365.25 days or 327.85 days for a dasha year'', and its
screenshot shows 365.25 selected \tagDOC\ \parencite[pp.~193--194]{geovision2006pl7}. The
code says something different \tagBIN. The OK handler (\texttt{0x4c7dc2}) writes setting byte
\texttt{0x68c}: 2 for the ``360'' button, 1 for ``365.25'' and 3 for ``327.85''; the
reset-to-defaults routine also sets~1. The year-length function \texttt{fcn.00464743} then
returns 360.0 for~2 and 327.85 for~3, but for~1 it returns \emph{365.2421875} days (the
tropical year), not 365.25. The Vimshottari start is computed as birth minus year length
times balance (\texttt{fcn.00460018}), and the year-length function is called from six places in
the dasha routines (\texttt{0x4600e8}, \texttt{0x4601af}, \texttt{0x460600}, \texttt{0x460697},
\texttt{0x4606b9}, \texttt{0x4709b4}). The constant 365.25 does occur in the same code region, but
next to 30.6001, 122.1, 1524 and 1\,720\,994.5: the Meeus calendar-to-Julian-Day conversion, not
a dasha year. Your PL7 dasha screenshots used for the 40/40 check (RA Job, SG, MANTA) show
periods within about two years of the chart, consistent with the explanation below.

\begin{derivbox}
\textbf{Why the 40/40 validation could not see this.} The two year lengths differ by
$0.0078$~day, or $11.25$~minutes per year of dasha time. Your validation compared dates
near the chart moment, within about $\pm 5$ years, so the boundaries moved by under an hour
and fell on the same calendar date. For a natal chart the drift grows: about $1.9$~h at 10
years from birth and about $11$~h at 60 years, which is enough to move pratyantar and
sookshma boundaries by a day, and prana boundaries much more.
\end{derivbox}

For horary use, where the periods of interest lie within weeks of the chart, the two
programs agree. If exact PL7 parity is ever wanted for long-range natal dashas, Kala Chakra
would need 365.2421875-day years. Either way, the knowledge base's description of the PL7
convention as ``365.25-day years'' should be corrected.

% =============================================================================
\section{The KP section: calculation and display}
\label{sec:kp}
% =============================================================================

\textbf{Cusps.} Placidus, with the cusp as the beginning of the house \tagBIN\ \tagDOC\
\parencite[pp.~56, 200]{geovision2006pl7}.

\textbf{Ayanamsa.} The dictionary says KP printouts use the Krishnamurti ayanamsa
``regardless of the default set'' \parencite[pp.~218--219]{geovision2006pl7}. In the code,
the only places that write the values 7 or 8 to the ayanamsa selector are in the options
dialog. Two printing/model routines (\texttt{fcn.00553970}, \texttt{fcn.00578e80}) load the
selector from a stored record (field \texttt{+0x260}), which matches the manual's
per-astrologer ayanamsa setting for model printing \parencite[p.~165]{geovision2006pl7}.
Which KP ayanamsa (selector 7 or 8) a KP printout uses was \emph{not determined} \tagINF.

\textbf{Lords.} The planetary and cusp tables show the nakshatra lord and sub lord
\parencite[p.~145]{geovision2006pl7}. Sub-sub lords are not documented and no sub-sub
column exists \tagDOC.

\textbf{Display.} PL7 has a circular KP chart, whose ``inner ring shows the Rashis, the next
ring the Nakshatras, the next the houses'' \parencite[p.~200]{geovision2006pl7}, with a table
of sub lords, occupants, ``tenants'' (planets in the nakshatras of occupants), owner and
tenants of the owner. There are also four KP worksheets (Krishnamurti, Krishnamurti~II,
Krishnamurti Prashna, Krishnamurti Prashna~II) and two reports. ``KP -- Chart'' gives the
birth chart, cuspal chart and ``Bhava Details (Placidus System)''. ``KP -- Significators''
gives the houses signified by planets, graded from ``very strong significator'' down; ruling
planets (day lord, lagna sign/star/sub lord, Moon sign/star/sub lord); and Fortuna \tagDOC.

\textbf{Horary.} ``Krishnamurti Prashna'' is cast for the current time ``and the ascendant
is based on a number selected by the user'' (worksheet help). Two integer-input calls in the
binary are bounded at 249 (\texttt{0x553f5a}, \texttt{0x557a21}), so the number is almost
certainly 1--249 \tagINF. The number-to-ascendant table was not decoded.

\begin{keybox}\kidea For KP horary, Kala Chakra is ahead of PL7 in four concrete ways. It
gives a five-level sub-lord chain (PL7 stops at the sub). It computes the family's
cusp-number vectors and power calculation, which PL7 has no notion of. It replicates Solar
Fire to $\le 1''$. And it is not affected by the Moon-vector planet error of
\autoref{sec:moonvector}. PL7 is ahead in breadth: significator grading, ruling planets,
Fortuna, and 1--249 horary ascendants.\end{keybox}

% =============================================================================
\section{Jaimini: karakas, padas and special lagnas}
\label{sec:jaimini}
% =============================================================================

Jaimini is implemented in PL7 \tagDOC. Everything below comes from the documentation; none
of it exists in Kala Chakra.

\textbf{Chara karakas.} A calculation option chooses 8 karakas including Rahu or 7 without
\parencite[p.~192]{geovision2006pl7}. Separately, the worksheet offers ``Karakas(8)'' and
``Karakas(7)'', where the 7-karaka scheme merges the mother and children karakas. Rahu's
degree is measured backwards (30 minus its degree in the sign). The default is not
documented.

\textbf{Aspects.} Rashi drishti: ``fixed signs aspect all the movable signs except the one
next to it'', and vice versa; dual signs aspect each other
\parencite[pp.~216--217]{geovision2006pl7}.

\textbf{Arudha padas.} A ``Jaimini padas and lagnas'' worksheet shows the padas of all twelve
houses for the birth chart, Navamsha and Dashamsha. The rule is to count from the house to
its lord, then as many houses again \parencite[p.~222]{geovision2006pl7}. The exceptions are
stated two different ways:

\begin{itemize}[leftmargin=*]
  \item The tutorial (\texttt{pib11755}): lord in the 4th or 7th gives the 4th house as
  Arudha Lagna, and lord in the 1st gives the 10th.
  \item The bundled BPHS (\texttt{CHAP29}): a pada falling in the same house goes to the
  10th from it, and one falling in the 7th goes to the 4th from the original house.
\end{itemize}

\textbf{The two disagree when the lord is in the 7th}, and neither says which the engine
applies \tagDOC.

\textbf{Upapada.} ``Generally considered as the pada of the 12th house, though there are
different views'' \parencite[p.~230]{geovision2006pl7}.

\textbf{Karakamsha and Swamsha.} Both are offered; the Karakamsha chart ``takes the sign of
the Atma Karaka (AK) in the Navamsha as the ascendant'', drawn on the birth chart or on the
Navamsha \parencite[p.~57]{geovision2006pl7}.

\textbf{Special lagnas.} Bhava lagna advances $15^\circ$ per hour from the Sun's position at
sunrise, Hora lagna one sign per hour, and Ghatika lagna one sign per 24 minutes; Varnada,
Indu and Pranapada are also given. Sunrise is taken at the Sun's edge or its centre, as an
option \parencite[p.~192]{geovision2006pl7}; the default and the refraction handling are not
documented.

\textbf{Argala.} Planets in the 2nd, 4th, 11th and 5th from a point cause argala, obstructed
from the 12th, 10th, 3rd and 9th; the direction is reversed for Rahu and Ketu.

% =============================================================================
\section{Guna matching (compatibility)}
\label{sec:guna}
% =============================================================================

PL7 offers both systems \tagDOC\ \parencite[pp.~104--106]{geovision2006pl7}:

\begin{itemize}[leftmargin=*]
  \item \textbf{Ashtakoota} (North Indian), on the Moon's nakshatra: varna, vashya, tara,
  yoni, graha maitri, gana, bhakoot and nadi, totalling 36. Gana, bhakoot and nadi are
  flagged as \emph{maha doshas}, and the Pancha-nadi exception is mentioned. The points per
  koota are not stated in any readable line.
  \item \textbf{Dashakoota} (South Indian, after Kalaprakasika), with three totals printed
  (Kalaprakasika, B.~V.~Raman and Ashtakoota methods): ``a number below 18 is questionable,
  between 18 and 24 is good and above 24 is excellent''.
  \item \textbf{Manglik (Kuja) dosha}: Mars in the 1st, 2nd, 4th, 7th, 8th or 12th from the
  ascendant and from the Moon, with seven printed cancellation rules.
  \item Matching can also be run on planets other than the Moon, and a rashi-distance
  comparison is offered.
\end{itemize}

There are no options to configure the koota rules. Kala Chakra has no matching module.

\Needspace{8\baselineskip}
% =============================================================================
\section{The planetary ``summariser'' table}
\label{sec:summary}
% =============================================================================

The planetary-details table has fifteen selectable columns \tagDOC\
\parencite[p.~145]{geovision2006pl7}: degrees; retrograde/combust flag; sign; nakshatra;
pada; nakshatra lord/sub; dignity; shadbala; vimshopaka bala (10 vargas); varga
classification; avasthas in the 3-, 5-, 9- and 12-state schemes; and lordships. The cusp
table has degrees, sign, nakshatra, pada and nakshatra lord/sub.

Kala Chakra's tables show degree, star lord and the five-level sub chain, house placement,
ownership numbers and the power string. They have no retrograde flag (the engine requests
no planetary speed), no dignity, no shadbala, no vimshopaka and no avasthas.
